\clearpage\thispagestyle{empty}
\pdfbookmark[0]{Quellen zu Instrumente (8. August 1818)}{instruments1818aug-sources}
\begin{center}\textbf{Quellen zu Instrumente (8. August 1818)}\end{center}
\noindent Diese KI-Übertragung ins heutige Deutsch folgt dem Bericht \textit{Instrumente} in den \href{https://gdz.sub.uni-goettingen.de/id/PPN236007467}{\textit{Werken}, Band~VI} (Göttingen, 1874), S.~410--414. Der Bericht erschien zuerst unter der Überschrift Göttingen in den \textit{Göttingischen gelehrten Anzeigen}, Stück~127, 8. August 1818, S.~\href{https://books.google.com/books?id=Uy1ZAAAAcAAJ&pg=PA1257}{1257}--\href{https://books.google.com/books?id=Uy1ZAAAAcAAJ&pg=PA1267}{1267}. Die Links zum Erstabdruck führen zum Digitalisat in Google Books. Der Bericht verwendet die erste Person Plural und spricht von Gauß in der dritten Person.
\par\smallskip\noindent In der Tafel der Polarsternbeobachtungen und dem anschließenden Satz steht im Erstabdruck auf S.~\href{https://books.google.com/books?id=Uy1ZAAAAcAAJ&pg=PA1263}{1263} bei der Korrektur für den 26. Juni ein Dezimalkomma. Der Mittelwert der Korrekturen endet dort auf 83. In der Werkausgabe fehlt in dieser Tabellenzelle der Dezimalpunkt und der Mittelwert endet auf 88. Auch bei drei Uranuswerten auf S.~\href{https://books.google.com/books?id=Uy1ZAAAAcAAJ&pg=PA1264}{1264}--\href{https://books.google.com/books?id=Uy1ZAAAAcAAJ&pg=PA1265}{1265} hat der Erstabdruck Dezimalkommas, während in der Werkausgabe keine Dezimalpunkte sichtbar sind. Die Übertragung übernimmt die Lesarten der Werkausgabe. KI-Übertragung und KI-Satz.
\clearpage
